% !TeX root = ../collection-solutions.tex

\titledquestion{Market structure, entrepreneurship and inequality}

\citet{deb2024market} documents three US trends since the late 1980s: the entrepreneurship rate \emph{fell} (from about $11\%$ to $8\%$), markups \emph{rose}, and the top-1\% income share \emph{rose}.

\begin{parts}
    \part[4] What is the key takeaway of the paper? Explain how the author arrives at this conclusion: what kind of model is used, how is it taken to the data, and how is the contribution of each force quantified?

    \begin{solutionorbox}[12cm]
        \emph{Takeaway} (2 points): One force — the growing dominance of highly productive ``superstar'' entrepreneurs and firms, who deter the entry of other productive agents into their markets — jointly explains all three trends. Falling entrepreneurship, rising markups and rising top income inequality are three faces of the same phenomenon, not separate puzzles.

        \emph{How he gets there} (2 points):
        \begin{itemize}
            \item A general-equilibrium occupational-choice model (\`a la Lucas) in which a finite number of agents per market choose between being a worker and being an entrepreneur, and entrepreneurs compete \emph{strategically} (Cournot) with each other and with a large corporate firm.
            \item Estimated by simulated method of moments: six structural parameters are re-estimated for every survey wave 1988--2018 to match six moments from the SCF (entrepreneurship rate, top-1\% income share, share of entrepreneurs in the top 1\%) and from Compustat (markups, fixed-cost share).
            \item Counterfactual decomposition: moving one parameter at a time from its 1988 to its 2018 value. Rising fixed costs explain about $90\%$ of the entrepreneurship decline; the rising productivity premium of large corporate firms and falling number of agents per market each explain about half of the markup rise; rising productivity dispersion explains most of the top-1\% increase.
        \end{itemize}
        Grading: 1 point for the one-force-explains-three-trends takeaway, 1 point for entry deterrence by dominant firms; 1 point for the model + estimation strategy (SMM, SCF + Compustat), 1 point for the counterfactual logic with at least one magnitude.
    \end{solutionorbox}

    \part[4] A common intuition says: if markups and top incomes are rising, running a firm has become more lucrative — so \emph{more} people should become entrepreneurs. Explain intuitively why the opposite happens in this model. How does strategic competition break the entry logic of the standard (Lucas-type) occupational choice model?

    \begin{solutionorbox}[12cm]
        \begin{itemize}
            \item (1 point) In the standard Lucas model, the entry decision depends only on one's \emph{own} productivity and aggregate prices: there is a single economy-wide threshold, and anyone above it becomes an entrepreneur. More lucrative entrepreneurship then indeed attracts entry.
            \item (2 points) With strategic (Cournot) competition in finite markets, an agent's profit depends on \emph{who else} is in the market: it is own productivity \emph{relative to} competitors' that determines the achievable market share and markup. A productive agent facing a much more productive rival sees their prospective market share driven toward zero, so staying a worker becomes the best response — even though the same agent would profitably enter a weaker market. Dominant entrepreneurs thus \emph{deter} the entry of other productive agents: the stronger a market's superstar, the more productive the best agent who chooses to remain a worker.
            \item (1 point) As technology becomes more dispersed and fixed costs rise, entry falls; the surviving dominant few hold larger market shares, charge higher markups, and capture more income. Falling entrepreneurship, rising markups and a rising top-1\% share therefore emerge \emph{together} — and part of workers' productivity gains is absorbed by markups into entrepreneurial profits, amplifying top inequality.
        \end{itemize}
    \end{solutionorbox}
\end{parts}
